% !TeX program = xelatex
% !TeX encoding = UTF-8
% D1 Stage 3 (Revised per Review): Complete Rigorous Proof of Computational Hardness for General B_graph
% Compile: xelatex D1_Bgraph_NP_hardness_complete_proof_report.tex (twice)
\documentclass[11pt,a4paper]{ctexart}
\usepackage[margin=24mm,top=25mm,bottom=24mm]{geometry}
\usepackage{amsmath,amssymb,amsthm,mathtools,bm}
\usepackage{booktabs,tabularx,array,longtable}
\usepackage{enumitem}
\usepackage{microtype}
\usepackage{xcolor}
\usepackage{hyperref}
\usepackage{fancyhdr,lastpage}
\usepackage[most]{tcolorbox}
\usepackage{listings}

\setlength{\parindent}{2em}
\setlength{\parskip}{0.25em}
\setlist{itemsep=.18em,topsep=.3em}

\definecolor{ink}{HTML}{173047}
\definecolor{accent}{HTML}{156D80}
\definecolor{soft}{HTML}{F0F6F8}
\definecolor{codebg}{HTML}{F7FAFC}
\definecolor{codeframe}{HTML}{CBD5E1}
\definecolor{codekw}{HTML}{0F4C81}
\definecolor{codecm}{HTML}{4A5568}
\definecolor{codestr}{HTML}{2B6CB0}

\hypersetup{colorlinks=true,linkcolor=ink,urlcolor=accent,citecolor=accent}

\pagestyle{fancy}
\fancyhf{}
\fancyhead[L]{\small\textcolor{ink}{D1：一般图证书 $B_{\rm graph}$ 计算困难性的完整严格证明（审查修订版）}}
\fancyhead[R]{\small\textcolor{ink}{理论研究报告（第三阶段定稿）}}
\fancyfoot[C]{\small 第 \thepage\ 页 / 共 \pageref{LastPage} 页}
\renewcommand{\headrulewidth}{0.3pt}

\newtheorem{theorem}{定理}[section]
\newtheorem{lemma}[theorem]{引理}
\newtheorem{proposition}[theorem]{命题}
\newtheorem{corollary}[theorem]{推论}
\theoremstyle{definition}
\newtheorem{definition}[theorem]{定义}
\newtheorem{remark}[theorem]{说明}
\newenvironment{proofcn}{\par\noindent\textbf{证明.}\ }{\hfill$\square$\par}

\newtcolorbox{resultbox}{breakable,colback=soft,colframe=accent,boxrule=.7pt,arc=1.3mm,left=3mm,right=3mm,top=2mm,bottom=2mm}

\lstdefinestyle{pycode}{
    language=Python,
    basicstyle=\ttfamily\footnotesize,
    backgroundcolor=\color{codebg},
    frame=single,
    rulecolor=\color{codeframe},
    keywordstyle=\color{codekw}\bfseries,
    commentstyle=\color{codecm}\itshape,
    stringstyle=\color{codestr},
    numbers=left,
    numberstyle=\tiny\color{codecm},
    stepnumber=1,
    numbersep=6pt,
    breaklines=true,
    breakatwhitespace=true,
    showstringspaces=false,
    tabsize=4,
    captionpos=b
}

\newcommand{\GraphCert}{\mathsf{GraphCert}}
\newcommand{\Odd}{\mathcal C_{\rm odd}}
\newcommand{\Ind}{\mathcal I}
\newcommand{\RR}{\mathbb R}
\newcommand{\QQ}{\mathbb Q}
\newcommand{\FF}{\mathbb F}
\newcommand{\1}{\boldsymbol 1}
\newcommand{\0}{\boldsymbol 0}
\newcommand{\bgraph}{B_{\rm graph}}
\newcommand{\buncut}{B_{\rm uncut}}
\newcommand{\ct}{c^*}
\newcommand{\chif}{\chi_f}

\begin{document}

\begin{center}
{\color{accent}\rule{\linewidth}{1.4pt}}\\[8mm]
{\Huge\bfseries D1：一般图证书 $B_{\rm graph}$ 计算困难性的\\[3mm]完整严格数学证明（审查修订版）}\\[5mm]
{\Large\bfseries 锚定星—顶图、分数色数精确阈值定理与 9 分类器显式反例}\\[6mm]
{\large 2026 年 10 月 13 日\quad|\quad 稳健加权投票项目理论定稿}
\end{center}

\vspace{3mm}
\begin{resultbox}
\textbf{摘要.} 在依赖不确定下的稳健加权投票中，高精度区域 $R_0=\{\sum_i q_i\le 1\}$ 内的精确 minimax 错误率由独裁规则与最优图证书 $B_H(b)=\inf_{y\ge 0:\,g_H(y)>0} b^\top y / g_H(y)$ 共同决定，其中 $g_H(y)=\min_{S\subseteq V}\max\{y(E(S)),y(E(V\setminus S))\}$ 是加权明智二部划分值。本报告（已根据审查意见修订完善全部细节）给出一般 $B_H(b)$ 计算困难性的\textbf{完整、自包含数学证明}及配套精确复核代码：
\begin{enumerate}[label=(\arabic*)]
\item \textbf{9 顶点星—团图 $J_4^\circ$ 的全局精确值与 $R_0$ 内 9 分类器显式反例}：在不假定任何对称性的前提下，利用强度为 $2$ 的四状态正交阵列给出初等对偶下界证明，确立 $\boxed{B_{J_4}=B_{J_4^\circ}=14/5<3=c^*}$；由此直接构造出落在 $R_0$ 内的 9 分类器显式投票实例（$q_i=1/9$），其精确最优错误率为 $R_{\rm det}^*=7/90$，严格小于最佳单模型与最便宜奇环给出的 $\min\{q_{\min},c^*\}=1/12=7.5/90$；
\item \textbf{单中心锚定定理与锚定星—顶图 $H(G)$ 的精确代数理论}：证明通过共享唯一中心顶点 $B_0$ 并对零成本垂直匹配边赋予大证书权重，可将任意多个构件的内部保障量严格锁定在同侧 $A_S$，彻底修复旧方案中 $(1,0)+(0,1)=(1,1)$ 的正负侧逃逸漏洞；同时导出两侧线性差恒等式与互补配对消去定理；
\item \textbf{核心主定理（锚定星—顶图上的分数色数精确阈值定理）}：对任意非空简单图 $G=(V_G,E_G)$ 与有理成本参数 $\beta\ge 1$，构造含 $2|V_G|+1$ 个顶点的图 $H(G)$，其最小奇环成本恒为 $c_{H(G)}^*(1,\beta)=2+\beta$，且成立精确充要条件
\[
 \boxed{B_{H(G)}(1,\beta) < 2+\beta \iff \chif(G) > 2+\beta,\qquad B_{H(G)}(1,\beta) = 2+\beta \iff \chif(G) \le 2+\beta.}
\]
结合确定性的“色数 $\chi(G)\le k$ 对分数色数 $\chif(G)\ge k'$”间隙归约（Lund--Yannakakis 1994；Zuckerman 2007），严格证明：\textbf{精确计算 $B_H(b)$（或以 $1/\operatorname{poly}(n)$ 量级的加性误差近似）是强 NP-hard 的；判定奇环界是否紧（$B_H(b)=c_H^*(b)$）及判定 $B_H(b)\ge\tau$ 是 NP-hard 的；决策语言 $\GraphCert=\{(H,b,\tau):B_H(b)\le\tau\}$ 属于 $\Sigma_2^P$ 且是 coNP-hard 的}；上述结论在完全图 $K_K$ 的 $R_0$ 可行域内（以平凡上界 $b_{ij}=q$ 填充非支撑边、以错误互斥分布保证 $\mathcal P(q,b)\neq\varnothing$）同样成立；
\item \textbf{五分量奇-$K_5$ 提升的 $\FF_2$ 上同调完整分类}：证明在全部 $2^{10}=1024$ 种五分量奇-$K_5$ 端口分裂中，仅有同构于 9 顶点图 $J_4^\circ$（$K=9$）的 10 种星—团划分满足 $B_H<3$，其余 1014 种（含全部 $K\le 8$ 的提升）恒有 $B_H=3=c_H^*$。
\end{enumerate}
\end{resultbox}

\vspace{2mm}
\begin{center}
\small
\begin{tabularx}{\linewidth}{@{}p{.31\linewidth}X p{.22\linewidth}@{}}
\toprule
问题或结构对象 & 本文严格证出的结论 & 复杂性/证明性质 \\
\midrule
二元证书 $y\in\{0,1\}^E$ & $g_H(y)\ge 1\iff \operatorname{supp}(y)$ 非二部；最优二元证书退化为 $c_H^*(b)$ & 奇环坍缩引理（§1.2） \\
9 顶点构件 $J_4^\circ$ 与 $R_0$ 显式反例 & $\boxed{B_{J_4^\circ}=14/5<3}$；$K=9$ 投票实例 $R_{\rm det}^*=7/90<1/12$ & 4 状态正交阵列对偶（§2） \\
多构件共享中心 $B_0$ & 消除 $(1,0)+(0,1)=(1,1)$ 翻转逃逸，$A_S$ 严格可加，$g_H=N$ & 单中心锚定定理（§3.1） \\
锚定星—顶图 $H(G)$ 恒等式 & $A_x(a,w)-D_x(w)\equiv W-\sum_{v}(d_w(v)-a_v)x_v$ & 精确代数恒等式（§3.2） \\
分数色数精确阈值定理 & $\boxed{B_{H(G)}(1,\beta)<2+\beta \iff \chif(G)>2+\beta}$（$\beta\ge 1$） & 10 行充要条件证明（§4.2） \\
精确计算 $B_H(b)$（及 $1/\operatorname{poly}$ 近似） & \textbf{强 NP-hard}（含稀疏图与 $R_0$ 内完全图 $K_K$ 实例） & 确定性 many-one 归约（§4.3） \\
紧性判定 $B_H(b)=c_H^*(b)$ 及 $B_H\ge\tau$ & \textbf{NP-hard}（由 $\chi(G)\le k$ 对 $\chif(G)\ge k'$ 归约） & 确定性 many-one 归约（§4.3） \\
决策语言 $\GraphCert$（$B_H(b)\le\tau$） & 属于 $\Sigma_2^P$ 且为 \textbf{coNP-hard}；单分支区域必属 $\mathrm{coNP}$ & 对偶见证障碍定理（§4.4） \\
五分量奇-$K_5$ 提升分类 & 1024 种二分中仅 9 顶点 $J_4^\circ$ 的 10 个同构副本满足 $B_H<3$ & $\FF_2$ 上同调与穷举（§5） \\
\bottomrule
\end{tabularx}
\end{center}

\newpage
\tableofcontents
\newpage

\section{精确定义、齐次规范化与二元证书坍缩引理}

\subsection{问题设定与基本记号}
设 $H=(V,E)$ 为有限简单无向图，$b=(b_e)_{e\in E}\in\QQ_{\ge 0}^E$ 为给定的非负有理边成本。对于任意非负证书边权向量 $y=(y_e)_{e\in E}\in\RR_{\ge 0}^E$ 及顶点子集 $S\subseteq V$，定义两侧内部边权与明智二部划分（judicious bipartition）值：
\begin{equation}\label{eq:ADg}
 A_S(y)=\sum_{e\in E(S)}y_e,\qquad
 D_S(y)=\sum_{e\in E(V\setminus S)}y_e,\qquad
 g_H(y)=\min_{S\subseteq V}\max\{A_S(y),D_S(y)\},
\end{equation}
其中 $E(S)=\{uv\in E:u\in S,\ v\in S\}$。最优图证书比率定义为
\begin{equation}\label{eq:B_def}
 \boxed{B_H(b)=\inf_{y\in\RR_{\ge 0}^E:\,g_H(y)>0}\frac{b^\top y}{g_H(y)},}
\end{equation}
若 $H$ 为二部图（此时对一切 $y\ge 0$ 均有 $g_H(y)=0$），约定 $B_H(b)=+\infty$。记 $H$ 中全部简单奇环的集合为 $\Odd(H)$，最便宜奇环上界为
\begin{equation}\label{eq:cstar_def}
 c_H^*(b)=\min_{C\in\Odd(H)}\sum_{e\in C}b_e,\qquad B_H(b)\le c_H^*(b).
\end{equation}
在完全图 $H=K_K$ 上简记 $\bgraph(b)=B_{K_K}(b)$。

\begin{lemma}[紧规范化与 $L_1$ 成本扰动稳定性]\label{lem:lip}
对任意非二部有限图 $H=(V,E)$ 与非负成本向量 $b,b'\in\RR_{\ge 0}^E$，有
\begin{equation}\label{eq:compact_norm}
 B_H(b)=\min\bigl\{b^\top y:\ y\in[0,1]^E,\ g_H(y)\ge 1\bigr\},
\end{equation}
且最优值可达到。进一步，函数 $b\mapsto B_H(b)$ 在 $\RR_{\ge 0}^E$ 上满足全局 $L_1$-Lipschitz 连续性：
\begin{equation}\label{eq:l1_lip}
 \boxed{\bigl|B_H(b)-B_H(b')\bigr|\le \|b-b'\|_1=\sum_{e\in E}|b_e-b'_e|.}
\end{equation}
\end{lemma}
\begin{proofcn}
由一阶齐次性 $g_H(ty)=t g_H(y)$（$t>0$），可将分母归一化为 $g_H(y)\ge 1$。对任意满足 $g_H(y)\ge 1$ 的 $y\ge 0$，作逐坐标截断 $\hat y_e=\min\{1,y_e\}\in[0,1]$。对任意划分 $S\subseteq V$，若原内部和 $A_S(y)\ge 1$，则当 $E(S)$ 中所有边均满足 $y_e\le 1$ 时有 $A_S(\hat y)=A_S(y)\ge 1$；当 $E(S)$ 中存在某条边 $e$ 满足 $y_e>1$ 时，有 $A_S(\hat y)\ge \hat y_e=1$。对 $D_S$ 同理，故 $g_H(\hat y)\ge 1$ 且 $b^\top\hat y\le b^\top y$。紧集 $[0,1]^E$ 上有限个闭半空间并集的交仍为非空紧集，故最小值在某点 $y^*\in[0,1]^E$ 达到，式 \eqref{eq:compact_norm} 成立。

设 $y^*\in[0,1]^E$ 为 $B_H(b')$ 的最优解（满足 $g_H(y^*)\ge 1$ 且 $(b')^\top y^*=B_H(b')$），则 $y^*$ 对 $B_H(b)$ 亦可行，故
\[
 B_H(b)-B_H(b')\le b^\top y^*-(b')^\top y^*=\sum_{e\in E}(b_e-b'_e)y^*_e\le \sum_{e\in E}|b_e-b'_e|=\|b-b'\|_1.
\]
交换 $b$ 与 $b'$ 即得式 \eqref{eq:l1_lip}。
\end{proofcn}

\subsection{为什么整数 $0$--$1$ 证书无法编码 NP-hardness：奇环坍缩引理}
\begin{proposition}[二元证书的奇环坍缩引理]\label{prop:binary_collapse}
设 $y\in\{0,1\}^E$ 为任意二元边权向量，记支撑边集 $F=\operatorname{supp}(y)=\{e\in E:y_e=1\}$。则
\begin{equation}\label{eq:binary_iff}
 g_H(y)\ge 1 \quad\Longleftrightarrow\quad (V,F)\text{ 不是二部图（即 }F\text{ 至少包含一条简单奇环）}.
\end{equation}
因此，在二元约束 $y\in\{0,1\}^E$ 下的最小化值恒等于最便宜奇环成本：
\begin{equation}\label{eq:binary_cstar}
 \min\bigl\{b^\top y:\ y\in\{0,1\}^E,\ g_H(y)\ge 1\bigr\}=c_H^*(b),
\end{equation}
该值可在 $O(|V|(|E|+|V|\log|V|))$ 多项式时间内通过奇偶分层最短路精确求得。
\end{proposition}
\begin{proofcn}
由于 $y_e\in\{0,1\}$，对任意划分 $S\subseteq V$，$A_S(y)=|F\cap E(S)|$ 与 $D_S(y)=|F\cap E(V\setminus S)|$ 均为非负整数。因此
\[
 \max\{A_S(y),D_S(y)\}\ge 1 \quad\Longleftrightarrow\quad F\cap\bigl(E(S)\cup E(V\setminus S)\bigr)\neq\varnothing,
\]
即划分 $(S,V\setminus S)$ 不是子图 $(V,F)$ 的合法二染色。对所有 $S\subseteq V$ 同时成立该式，当且仅当 $(V,F)$ 不存在合法二染色，即 $(V,F)$ 非二部（含奇环 $C\in\Odd(H)$）。由于 $b\ge 0$，包含奇环的极小边集恰为单条简单奇环 $C\in\Odd(H)$，故式 \eqref{eq:binary_cstar} 成立。
\end{proofcn}

\section{星—团构件 $J_4^\circ$ 的全局最优性与 $R_0$ 内 9 分类器显式反例}

\subsection{从 10 顶点 $J_4$ 到 9 顶点约化图 $J_4^\circ$ 及无对称性对偶证明（完成任务 A）}
交接备忘录第 3.1 节引入了 10 顶点图 $J_4$：其中顶点 $T_0$ 度数为 $1$（仅与 $B_0$ 连有一条零成本悬挂桥边，不关联任何付费边）。删去度一顶点 $T_0$ 后得到 \textbf{9 顶点图 $J_4^\circ$}（顶点集 $\{B_0,T_1,T_2,T_3,T_4,B_1,B_2,B_3,B_4\}$，含 4 条零成本垂直匹配边 $T_iB_i$、底层 4 条成本为 $1$ 的星边 $B_0B_i$ 与顶层 6 条成本为 $1$ 的团边 $T_iT_j$），满足 $B_{J_4^\circ}(b)=B_{J_4}(b)$。

\begin{theorem}[星—团构件 $J_4$ 与 $J_4^\circ$ 的精确全局最优值]\label{thm:J4_exact}
对于 10 顶点图 $J_4$ 与 9 顶点图 $J_4^\circ$，其最便宜奇环成本为 $c_{J_4}^*(b)=c_{J_4^\circ}^*(b)=3$，而最优图证书的精确全局最优值为
\begin{equation}\label{eq:J4_val}
 \boxed{B_{J_4}(b)=B_{J_4^\circ}(b)=\frac{14}{5}=2.8<3.}
\end{equation}
\end{theorem}
\begin{proofcn}
\textbf{第一步（上界与奇环成本）.}
顶层 $K_4$ 内部最短奇环为三角形 $(T_i,T_j,T_k)$，付费成本为 $3$；经过中心 $B_0$ 的任何奇环必含两条星边 $B_0B_i,B_0B_j$、两条零成本垂直边 $B_iT_i,B_jT_j$ 以及顶层奇数长路径（至少 1 条团边），付费成本至少为 $2+1=3$，故 $c_{J_4^\circ}^*(b)=3$。
取星边权重 $a_i=2/5$（$1\le i\le 4$）、团边权重 $w_{ij}=1/5$（$1\le i<j\le 4$）、垂直零成本边权重 $y_{T_iB_i}=1$，总付费成本为 $4\times\frac{2}{5}+6\times\frac{1}{5}=\frac{14}{5}$。固定 $B_0\in S$，若垂直边均异色且 $k=\sum_{i=1}^4 \1\{B_i\in S\}\in\{0,1,2,3,4\}$，则 $A_x=\frac{2}{5}k+\frac{1}{5}\binom{4-k}{2}\in\{\frac{6}{5},1,1,\frac{6}{5},\frac{8}{5}\}\ge 1$，故 $g_{J_4^\circ}(y)=1$，$B_{J_4^\circ}(b)\le 14/5$。

\medskip
\noindent\textbf{第二步（对任意非负证书 $(a,w)\ge 0$ 的全局下界 $\ge 14/5$）.}
设 $(a,w)\ge 0$ 为任意满足 $g_{J_4^\circ}(y)\ge 1$ 的可行证书。固定 $B_0\in S$，对每个 $x=(x_1,x_2,x_3,x_4)\in\{0,1\}^4$ 考察垂直异色划分 $S_x$（$B_i\in S\iff x_i=1\iff T_i\notin S$），有
\begin{equation}\label{eq:J4_AD}
 A_x(a,w)=\sum_{i=1}^4 a_i x_i + \sum_{1\le i<j\le 4} w_{ij}(1-x_i)(1-x_j),\qquad
 D_x(w)=\sum_{1\le i<j\le 4} w_{ij} x_i x_j.
\end{equation}
分两种互补情形讨论：
\begin{itemize}
\item \textbf{情形 1：所有团边均满足 $w_{ij}<1$（$1\le i<j\le 4$）.}
此时对任意满足 $|x|\in\{1,2\}$ 的 10 个状态 $x$，均有 $D_x(w)<1$，故这 10 个状态必须全部满足 $A_x(a,w)\ge 1$。将 $|x|=1$ 的 4 式乘以 $2/5$、$|x|=2$ 的 6 式乘以 $1/5$ 求和，因每个 $a_i$ 与 $w_{ij}$ 的组合系数均为 $\frac{2}{5}\cdot 1+\frac{1}{5}\cdot 3=1$ 与 $\frac{2}{5}\cdot 2+\frac{1}{5}\cdot 1=1$，得
\[
 \sum_{i=1}^4 a_i + \sum_{1\le i<j\le 4} w_{ij}
 = \frac{2}{5}\sum_{|x|=1} A_x(a,w) + \frac{1}{5}\sum_{|x|=2} A_x(a,w)
 \ge \frac{2}{5}\times 4 + \frac{1}{5}\times 6 = \frac{14}{5}.
\]
\item \textbf{情形 2：存在某条团边满足 $w_{uv}\ge 1$（不妨设 $w_{12}\ge 1$）.}
取 4 个正交阵列状态 $s^{(1)}=(0,1,0,0),\, s^{(2)}=(0,1,1,1),\, s^{(3)}=(1,0,0,1),\, s^{(4)}=(1,0,1,0)$。在全部 4 个状态中 $x_1\neq x_2$（故 $w_{12}$ 恒被割开），而其余每个变量在 $\sum_{r=1}^4(A_{s^{(r)}}+D_{s^{(r)}})$ 中恰好各出现 2 次，故
\[
 \sum_{i=1}^4 a_i + \sum_{1\le i<j\le 4} w_{ij}
 = w_{12} + \frac{1}{2}\sum_{r=1}^4\bigl(A_{s^{(r)}}(a,w)+D_{s^{(r)}}(w)\bigr)
 \ge 1 + \frac{1}{2}\times 4 = 3 > \frac{14}{5}.
\]
\end{itemize}
综上，$B_{J_4}(b)=B_{J_4^\circ}(b)=14/5$。
\end{proofcn}

\subsection{$R_0$ 内严格突破“最佳单模型与最便宜奇环”的 9 分类器显式实例}
将 9 顶点图 $J_4^\circ$ 直接嵌入 $K=9$ 的稳健加权投票模型，即可得到选题文档结果 4(b) 所需的显式小规模实例（无需随机图存在性论证）：

\begin{corollary}[$R_0$ 内的 9 分类器显式反例]\label{cor:K9_voting_example}
取 $K=9$ 个分类器，对应 $J_4^\circ$ 的顶点集 $V=\{B_0,T_1,T_2,T_3,T_4,B_1,B_2,B_3,B_4\}$。令每个分类器的边际错误率为 $q_i=\frac{1}{9}$（满足 $\sum_{i=1}^9 q_i=1$，位于 $R_0$ 内），成对共同出错概率上界 $b_{uv}$ 设定为：
\begin{equation}\label{eq:K9_b_def}
 b_{uv}=\begin{cases}
 0, & \text{若 } uv\in E_0=\{T_1B_1,T_2B_2,T_3B_3,T_4B_4\},\\[3pt]
 \dfrac{1}{36}, & \text{若 } uv\in E_{\rm star}\cup E_{\rm top}\text{（即 }J_4^\circ\text{ 的 10 条水平付费边）},\\[5pt]
 \dfrac{1}{9}=q, & \text{若 } uv\notin E(J_4^\circ)\text{（其余 22 条非支撑填充边，取平凡边际上界）}.
 \end{cases}
\end{equation}
则可行联合分布集 $\mathcal P(q,b)$ 非空，且在 $R_0$ 内有：
\begin{equation}\label{eq:K9_gap}
 \boxed{R_{\rm det}^*(q,b)=\bgraph(b)=\frac{7}{90}\approx 0.0778 \;<\; \frac{1}{12}=\frac{7.5}{90}=\min\{q_{\min},\,c^*(b)\}.}
\end{equation}
\end{corollary}
\begin{proofcn}
\textbf{1. 分布集 $\mathcal P(q,b)$ 的非空性（错误互斥分布）.}
由于 $\sum_{i=1}^9 q_i = 9\times\frac{1}{9}=1$，可在样本空间 $[0,1]$（配备 Lebesgue 测度）上令第 $i$ 个分类器的出错事件为互不相交的区间 $\bigl[\frac{i-1}{9},\frac{i}{9}\bigr)$（$i=1,\dots,9$）。在该“错误互斥”联合分布 $\pi$ 下，$\Pr_\pi(X_i=1)=\frac{1}{9}=q_i$，且对任意 $i\neq j$ 均有 $\Pr_\pi(X_i=1,X_j=1)=0\le b_{ij}$，故 $\pi\in\mathcal P(q,b)\neq\varnothing$。

\textbf{2. 精确计算 $c^*(b)$ 与 $\bgraph(b)$.}
在 $J_4^\circ$ 的支撑边上，最便宜奇环含 3 条成本为 $\frac{1}{36}$ 的边，成本为 $\frac{3}{36}=\frac{1}{12}$；任何包含至少一条非支撑填充边（成本 $\frac{1}{9}>\frac{1}{12}$）的奇环成本均 $>\frac{1}{12}$，故 $c^*(b)=\frac{1}{12}$，且 $\min\{q_{\min},c^*(b)\}=\min\{\frac{1}{9},\frac{1}{12}\}=\frac{1}{12}$。
另一方面，因所有 22 条非支撑填充边的成本 $\frac{1}{9}$ 严格大于已知上界 $U=c^*(b)=\frac{1}{12}$，由高成本边安全剔除定理（第一阶段报告定理 6.1，严格不等式 $b_e>U$ 成立），所有填充边可安全删去而不改变 $\bgraph(b)$。再由定理~\ref{thm:J4_exact} 并按比例因子 $\gamma=\frac{1}{36}$ 缩放，得
\[
 \bgraph(b)=\frac{1}{36}B_{J_4^\circ}(\1)=\frac{1}{36}\times\frac{14}{5}=\frac{7}{90}<\frac{1}{12},
\]
从而由 $R_0$ 内的主定理得 $R_{\rm det}^*(q,b)=\min\{q_{\min},\bgraph(b)\}=\frac{7}{90}$。
\end{proofcn}

\section{单中心锚定定理与锚定星—顶图 $H(G)$ 的代数结构}

\subsection{单中心锚定定理：修复跨构件正负侧逃逸（完成优先任务 B）}
\begin{theorem}[单中心锚定与跨构件严格可加性定理]\label{thm:anchor_additive}
设 $G_1=(V_1,E_1),\dots,G_N=(V_N,E_N)$ 为 $N$ 个顶层图，构造共享唯一中心顶点 $B_0$ 的合并星—顶图 $H$：对每个构件 $r\in\{1,\dots,N\}$ 及顶点 $v\in V_r$，设有一条零成本垂直边 $T_v^{(r)}B_v^{(r)}$（成本 $b=0$）、一条连接共享中心 $B_0$ 的底层星边 $B_0B_v^{(r)}$，以及顶层边 $\{T_u^{(r)}T_v^{(r)}:uv\in E_r\}$。
若在第 $r$ 个构件上给定证书边权 $(a^{(r)},w^{(r)})\ge 0$，使其在固定 $B_0\in S$ 的单侧内部函数满足 $\min_{x^{(r)}\in\{0,1\}^{V_r}} A_{x^{(r)}}^{(r)}(a^{(r)},w^{(r)})\ge \gamma_r > 0$（$r=1,\dots,N$），并将所有零成本垂直边的证书权重设为 $y_{\rm vert}=\Gamma:=\sum_{r=1}^N \gamma_r$，则在合并图 $H$ 上有精确加性下界：
\begin{equation}\label{eq:strict_add}
 \boxed{g_H(y)\ge \sum_{r=1}^N \gamma_r = \Gamma.}
\end{equation}
\end{theorem}
\begin{proofcn}
任取划分 $S\subseteq V(H)$，由互补对称性可设共享中心 $B_0\in S$。若某条零成本垂直边两端同色，则仅该边已提供内部权重 $\Gamma$；若全部垂直边均异色，则因 $B_0\in S$ 对所有构件同时固定在 $S$ 侧，各构件在 $S$ 侧的内部边权直接相加：$A_S(y)=\sum_{r=1}^N A_{x^{(r)}}^{(r)}(a^{(r)},w^{(r)})\ge \sum_{r=1}^N \gamma_r=\Gamma$。
\end{proofcn}

\subsection{一般锚定星—顶图 $H(G)$ 的两个基本代数恒等式}
\begin{definition}[锚定星—顶图 $H(G)$]\label{def:HG}
给定简单无向图 $G=(V_G,E_G)$（$n=|V_G|,\, m=|E_G|$），定义锚定星—顶图 $H(G)=(V,E_0\dot\cup E_{\rm star}\dot\cup E_{\rm top})$：
顶点集 $V=\{B_0\}\cup\{T_v,B_v:v\in V_G\}$（共 $2n+1$ 个顶点）；零成本垂直匹配边 $E_0=\{T_vB_v:v\in V_G\}$（成本 $0$）；底层星边 $E_{\rm star}=\{B_0B_v:v\in V_G\}$（成本 $b_v^{\rm star}$）；顶层图边 $E_{\rm top}=\{T_uT_v:uv\in E_G\}$（成本 $b_{uv}^{\rm top}$）。当 $b_v^{\rm star}=1,\, b_{uv}^{\rm top}=\beta$ 时简记 $B_{H(G)}(b)$ 为 $B_{H(G)}(1,\beta)$。
\end{definition}

固定 $B_0\in S$，记 $x_v=\1\{B_v\in S\}=1-\1\{T_v\in S\}\in\{0,1\}$，并记 $A=\sum_{v\in V_G}a_v,\, W=\sum_{uv\in E_G}w_{uv},\, d_w(v)=\sum_{u:\,uv\in E_G}w_{uv}$。对任意状态 $X=\{v\in V_G:x_v=1\}$，有
\begin{equation}\label{eq:HG_AD}
 A_X(a,w)=\sum_{v\in V_G}a_v x_v + \sum_{uv\in E_G}w_{uv}(1-x_u)(1-x_v),\qquad
 D_X(w)=\sum_{uv\in E_G}w_{uv}x_u x_v.
\end{equation}

\begin{theorem}[两侧线性差恒等式与互补配对消去定理]\label{thm:HG_identities}
在锚定星—顶图 $H(G)$ 上，对任意证书权重 $(a,w)\ge 0$：
\begin{enumerate}[label=(\alph*)]
\item \textbf{两侧线性差恒等式}：对每个状态 $x\in\{0,1\}^{V_G}$，有
\begin{equation}\label{eq:linear_diff}
 \boxed{A_x(a,w) - D_x(w) \equiv W - \sum_{v\in V_G}\bigl(d_w(v)-a_v\bigr)x_v.}
\end{equation}
特别地，对 $G$ 的任意独立集 $I\in\Ind(G)$，恒有 $D_I(w)=0$ 且 $A_I(a,w) = W - \sum_{v\in I}(d_w(v)-a_v)$。
\item \textbf{割配对星权消去与未割下界}：对任意可行证书 $(a,w)$，顶层未割值 $u_G(w)=W-\operatorname{MaxCut}(G,w)$ 满足
\begin{equation}\label{eq:uncut_duality}
 \boxed{u_G(w)\ge \min\{1,\; 2-A\}.}
\end{equation}
\item \textbf{独立集—点覆盖配对星权消去}：对任意可行证书 $(a,w)$，有
\begin{equation}\label{eq:alpha_deg_bound}
 \boxed{\alpha(G;\,d_w):=\max_{I\in\Ind(G)}\sum_{v\in I}d_w(v)\le \max\{W-1,\; A+W-2\}.}
\end{equation}
\end{enumerate}
\end{theorem}
\begin{proofcn}
(a) 由 $(1-x_u)(1-x_v)-x_u x_v = 1 - x_u - x_v$ 代入式 \eqref{eq:HG_AD} 即得式 \eqref{eq:linear_diff}；独立集 $I$ 内部无边故 $D_I(w)=0$。
(b) 对任意 $X\subseteq V_G$，记 $U_1=w(E(G[X])),\, U_0=w(E(G[V_G\setminus X]))$。若 $\max\{U_0,U_1\}\ge 1$，则 $U(w;X)=U_0+U_1\ge 1$；若 $\max\{U_0,U_1\}<1$，则 $D_X=U_1<1$ 且 $D_{V_G\setminus X}=U_0<1$，故 $A_X=a(X)+U_0\ge 1$ 且 $A_{V_G\setminus X}=a(V_G\setminus X)+U_1\ge 1$，相加得 $A+U(w;X)\ge 2$。
(c) 对任意 $I\in\Ind(G)$，状态 $I$ 给出 $W-\sum_{v\in I}d_w(v)+\sum_{v\in I}a_v\ge 1$；状态 $V_G\setminus I$ 给出 $A_{V_G\setminus I}=A-\sum_{v\in I}a_v$ 与 $D_{V_G\setminus I}=W-\sum_{v\in I}d_w(v)$，分 $D_{V_G\setminus I}\ge 1$ 与 $<1$ 两种情形相加消去 $\sum_{v\in I}a_v$ 即得式 \eqref{eq:alpha_deg_bound}。
\end{proofcn}

\section{主定理：分数色数精确阈值定理与一般 $B_{\rm graph}$ 的 NP-hardness}

\subsection{分数色数 $\chif(G)$ 的对偶定义与非孤立支撑引理}
\begin{definition}[分数色数及其对偶线性规划]\label{def:chif}
对有限简单图 $G=(V_G,E_G)$，令 $\Ind(G)$ 为 $G$ 的全体独立集构成的集合族。图 $G$ 的分数色数 $\chif(G)$ 由如下互为对偶的线性规划定义：
\begin{equation}\label{eq:chif_lp}
 \chif(G)
 = \min\left\{\sum_{I\in\Ind(G)}\lambda_I:\ \sum_{I\ni v}\lambda_I\ge 1\ (\forall v\in V_G),\ \lambda_I\ge 0\right\}
 = \max\left\{\sum_{v\in V_G}c_v:\ \sum_{v\in I}c_v\le 1\ (\forall I\in\Ind(G)),\ c_v\ge 0\right\}.
\end{equation}
对任意图 $G$，恒有 $\omega(G)\le \overline{\vartheta}(G)\le \chif(G)\le \chi(G)$，以及 $\chif(G)\ge |V_G|/\alpha(G)$。
\end{definition}

\begin{lemma}[非孤立顶点支撑性质]\label{lem:no_isolated}
设 $G=(V_G,E_G)$ 满足 $\chif(G)>1$。则存在式 \eqref{eq:chif_lp} 右端对偶问题的最优分数团向量 $c^*\in\QQ_{\ge 0}^{V_G}$，满足 $\sum_{v\in V_G}c_v^*=\chif(G)$，且对 $G$ 的一切孤立顶点 $v$（即 $\deg_G(v)=0$）均有 $c_v^*=0$。
\end{lemma}
\begin{proofcn}
若 $v_0\in V_G$ 为孤立顶点，则对 $G'=G\setminus\{v_0\}$ 的任意独立集 $I'\in\Ind(G')$，$I'\cup\{v_0\}$ 均为 $G$ 的独立集。由于 $\chif(G)>1$，$G$ 至少含有一条边，故 $\chif(G)=\chif(G')$。取 $G'$ 上对偶问题的最优有理解 $\tilde c\in\QQ_{\ge 0}^{V_G\setminus\{v_0\}}$，延拓为 $c_{v_0}^*=0$ 及 $c_v^*=\tilde c_v$（$v\neq v_0$），则对任意 $I\in\Ind(G)$ 有 $\sum_{v\in I}c_v^*\le 1$ 且 $\sum_{v\in V_G}c_v^*=\chif(G)$。归纳删去所有孤立顶点即证。
\end{proofcn}

\subsection{核心主定理：锚定星—顶图上的分数色数精确阈值定理（完成任务 C）}
\begin{theorem}[锚定星—顶图上的分数色数精确阈值定理]\label{thm:main_threshold}
设 $G=(V_G,E_G)$ 为任意至少含一条边的有限简单无向图（$m=|E_G|\ge 1,n=|V_G|$），且成本参数 $\beta\in\QQ$ 满足 $\beta\ge 1$（若 $G$ 无三角形，则只需 $\beta\ge 1/2$）。在锚定星—顶图 $H(G)$ 上赋予零成本垂直边 $b_e=0$、底层星边 $b_v^{\rm star}=1$、顶层图边 $b_{uv}^{\rm top}=\beta$。则：
\begin{enumerate}[label=(\arabic*)]
\item $H(G)$ 的最小奇环成本恒为 $c_{H(G)}^*(1,\beta)=2+\beta$；
\item 图证书严格突破奇环上界当且仅当 $G$ 的分数色数严格大于 $2+\beta$：
\begin{equation}\label{eq:main_iff}
 \boxed{B_{H(G)}(1,\beta) < 2+\beta \quad\Longleftrightarrow\quad \chif(G) > 2+\beta,}
\end{equation}
等价地，奇环上界精确紧（$B_{H(G)}(1,\beta)=c_{H(G)}^*(1,\beta)$）当且仅当 $\chif(G)\le 2+\beta$：
\begin{equation}\label{eq:main_eq}
 \boxed{B_{H(G)}(1,\beta) = 2+\beta \quad\Longleftrightarrow\quad \chif(G) \le 2+\beta;}
\end{equation}
\item \textbf{显式逆多项式间隙估计}：若 $\chif(G)\ge (2+\beta)+\delta$（其中 $\delta>0$），则取 $\varepsilon=\frac{1}{2m\chif(G)}\ge \frac{1}{2mn}$ 时有显式上界
\begin{equation}\label{eq:explicit_gap}
 \boxed{B_{H(G)}(1,\beta) \le (2+\beta) - \frac{\delta}{2m\chif(G)} \le (2+\beta) - \frac{\delta}{2mn}.}
\end{equation}
\end{enumerate}
\end{theorem}
\begin{proofcn}
\textbf{第一步：证明 $c_{H(G)}^*(1,\beta)=2+\beta$.}
任取 $uv\in E_G$，五元环 $(B_0-B_u-T_u-T_v-B_v-B_0)$ 含两条星边、一条顶层边与两条零成本垂直边，总长为 $5$（奇数），付费成本为 $2+\beta$。反之，经过 $B_0$ 的任何简单奇环成本至少为 $2+\beta$；完全位于顶层 $G$ 内的简单奇环长度至少为 $3$（无三角形时至少为 $5$），当 $\beta\ge 1$（或无三角形且 $\beta\ge 1/2$）时其成本 $\ge 3\beta\ge 2+\beta$（或 $5\beta\ge 2+\beta$）。故 $c_{H(G)}^*(1,\beta)=2+\beta$。

\medskip
\noindent\textbf{第二步：证明 $\chif(G)\le 2+\beta \implies B_{H(G)}(1,\beta)=2+\beta$.}
设 $(a,w)\ge 0$ 满足 $g_{H(G)}(a,w)\ge 1$。
在空集状态 $I=\varnothing$ 下，$D_\varnothing(w)=0<1$，故 $A_\varnothing(a,w)=W\ge 1$。
在任意独立集状态 $I\in\Ind(G)$ 下，$D_I(w)=0<1$，故 $A_I(a,w)=W-\sum_{v\in I}(d_w(v)-a_v)\ge 1$，即 $\sum_{v\in I}(d_w(v)-a_v)\le W-1$。
定义 $c_v:=\max\{0,\,d_w(v)-a_v\}\ge 0$。对任意 $I\in\Ind(G)$，子集 $I^+:=\{v\in I:d_w(v)-a_v>0\}\subseteq I$ 仍是独立集，故
\[
 \sum_{v\in I}c_v = \sum_{v\in I^+}\bigl(d_w(v)-a_v\bigr) \le W-1\qquad \forall I\in\Ind(G).
\]
由式 \eqref{eq:chif_lp} 的对偶定义知 $\sum_{v\in V_G}c_v \le (W-1)\chif(G)$。又因 $a_v\ge d_w(v)-c_v$，对 $v\in V_G$ 求和得 $A\ge 2W-\sum_{v\in V_G}c_v\ge 2W-(W-1)\chif(G)$。两端加 $\beta W$ 并利用 $W\ge 1$ 与 $\chif(G)\le 2+\beta$，得
\[
 A + \beta W \ge (2+\beta)W - (W-1)\chif(G) = (2+\beta) + (W-1)\bigl(2+\beta-\chif(G)\bigr) \ge 2+\beta.
\]

\medskip
\noindent\textbf{第三步：证明 $\chif(G)>2+\beta \implies B_{H(G)}(1,\beta)<2+\beta$ 及式 \eqref{eq:explicit_gap}.}
设 $\chif(G)=(2+\beta)+\delta$（$\delta>0$）。由引理~\ref{lem:no_isolated} 取对偶最优解 $c^*\in\QQ_{\ge 0}^{V_G}$（满足 $\sum_{v\in I}c_v^*\le 1$，$\sum_{v}c_v^*=\chif(G)$，且在孤立点上 $c_v^*=0$）。令 $\varepsilon=\frac{1}{2m\chif(G)}>0$，并设
\[
 w_{uv} := \frac{1+\varepsilon}{m}\quad (\forall uv\in E_G),\qquad
 a_v := d_w(v) - \varepsilon c_v^* = \frac{1+\varepsilon}{m}\deg_G(v) - \varepsilon c_v^*\quad (\forall v\in V_G).
\]
若 $\deg_G(v)=0$ 则 $a_v=0$；若 $\deg_G(v)\ge 1$ 则 $a_v\ge \frac{1+\varepsilon}{m}-\varepsilon > 0$。对任意 $X\subseteq V_G$，由式 \eqref{eq:linear_diff} 有
\[
 A_X(a,w) = 1 + \varepsilon\left(1 - \sum_{v\in X}c_v^*\right) + \sum_{uv\in E(G[X])}w_{uv}.
\]
若 $X\in\Ind(G)$，则 $\sum_{v\in X}c_v^*\le 1$，故 $A_X(a,w)\ge 1$；若 $X\notin\Ind(G)$，则 $G[X]$ 至少含一条边，故 $\sum_{uv\in E(G[X])}w_{uv}\ge \frac{1+\varepsilon}{m}$，从而 $A_X(a,w)\ge 1-\varepsilon\chif(G)+\frac{1+\varepsilon}{m}=1+\frac{1+2\varepsilon}{2m}>1$。总成本为
\[
 A + \beta W = (2+\beta)(1+\varepsilon) - \varepsilon\chif(G) = (2+\beta) - \frac{\delta}{2m\chif(G)} \le (2+\beta) - \frac{\delta}{2mn}.
\]
证毕。
\end{proofcn}

\begin{remark}[团数上界公式及其为何不给出常数近似困难性]\label{rem:lovasz_theta}
若图 $G$ 含有 $k$-团（$\omega(G)\ge k\ge 4$），在 $H(G)$ 中仅对该 $k$-团对应的子图 $J_k^\circ$ 赋权 $a_v=\frac{2}{k+1},\, w_{uv}=\frac{2}{(k-2)(k+1)}$，可得组合上界 $B_{H(G)}(1,k-3)\le (k-1)-\frac{2}{(k-2)(k+1)}$（例如 $k=4,\beta=1$ 时给出 $14/5=2.8$）。然而，正如审查意见所指出的，这一公式\textbf{不能}用于导出常数加性误差的 NP-hardness：对于固定常数 $k$，检测 $\omega(G)\ge k$ 可在 $O(n^k)$ 多项式时间内完成；对于一般的 $k$，Lov\'asz $\vartheta$ 函数满足夹逼定理 $\omega(G)\le\overline{\vartheta}(G)\le\chif(G)$ 且可通过半定规划在多项式时间内计算，故任何多项式时间归约都不可能在 YES/NO 实例间造成 $\omega(G)\ge k$ 与 $\chif(G)\le k-1$ 的分离（除非 $\mathbf{P}=\mathbf{NP}$）。因此，下述困难性定理严格建立在“$\chi(G)\le k$ 对 $\chif(G)\ge k'$”的间隙归约之上，所得近似困难性精度为逆多项式量级 $\Omega(1/\operatorname{poly}(n))$。
\end{remark}

\subsection{一般 $B_{\rm graph}$ 的计算困难性主定理（修正完善版）}
\begin{theorem}[一般 $B_{\rm graph}$ 的计算困难性主定理]\label{thm:main_hardness}
在标准有理数二进制编码下，成立如下确定性多项式时间归约与复杂性下界：
\begin{enumerate}[label=(\alph*)]
\item \textbf{奇环紧性判定与上界阈值的 NP-hardness}：判定给定实例 $(H,b)$ 是否满足 $B_H(b)=c_H^*(b)$，以及给定 $(H,b,\tau)$ 判定是否 $B_H(b)\ge \tau$，在确定性多项式时间 many-one（Karp）归约下均是 \textbf{NP-hard} 的。
\item \textbf{决策语言 $\GraphCert$ 的 coNP-hardness}：决策语言 $\GraphCert=\{(H,b,\tau):B_H(b)\le\tau\}$ 属于 $\Sigma_2^P$，且在确定性多项式时间 many-one（Karp）归约下是 \textbf{coNP-hard} 的。
\item \textbf{精确求值与逆多项式加性近似的 NP-hardness}：精确计算函数值 $B_H(b)$（或以 $\frac{1}{4mn}$ 的逆多项式加性误差近似计算 $B_H(b)$）是 \textbf{NP-hard} 的。
\item \textbf{稳健加权投票完全图 $K_K$（落在 $R_0$ 内部）上的困难性}：上述 (a)--(c) 的困难性结论在稳健加权投票原始完全图 $H=K_K$ 且参数 $(q,b)$ 严格满足 $\sum_{i=1}^K q_i\le 1$ 及 $\mathcal P(q,b)\neq\varnothing$ 的实例上同样成立。
\end{enumerate}
\end{theorem}
\begin{proofcn}
\textbf{1. 源问题的确定性间隙归约.}
由 Gr\"otschel--Lov\'asz--Schrijver~\cite{GLS1981,GLS1988}（定理 9.3.30），精确判定 $\chif(G)\le k$ 是 NP-complete 的；进一步，由 Lund--Yannakakis~\cite{LundYannakakis1994}（定理 2.5，利用 $\chif(G)\ge |V_G|/\alpha(G)$）及 Zuckerman~\cite{Zuckerman2007}（定理 1.2，通过线性度数提取器将 Feige--Kilian / H\aa stad 的团与色数不可近似性去随机化为确定性多项式时间 many-one 归约），存在确定性多项式时间归约，将任意 \textsc{3-SAT} 公式 $\Phi$ 映射为简单图 $G_\Phi=(V_G,E_G)$（$n=|V_G|,m=|E_G|\ge 1$）与多项式有界整数 $k\ge 3, k'\ge k+1$，使得：
\begin{itemize}
\item 若 $\Phi\in\textsc{3-SAT}$，则 $\chi(G_\Phi)\le k$，从而 $\chif(G_\Phi)\le \chi(G_\Phi)\le k$；
\item 若 $\Phi\notin\textsc{3-SAT}$，则 $\chif(G_\Phi)\ge |V_G|/\alpha(G_\Phi)\ge k'\ge k+1$。
\end{itemize}

\textbf{2. 证明 (a), (b), (c).}
对 $G=G_\Phi$ 构造定义~\ref{def:HG} 的锚定星—顶图 $H(G)$（顶点数 $K=2n+1$），设定垂直边成本为 $0$，星边成本为 $1$，顶层边成本为整数 $\beta:=k-2\ge 1$。由定理~\ref{thm:main_threshold}：
\begin{itemize}
\item 若 $\Phi\in\textsc{3-SAT}$，则 $\chif(G)\le k=2+\beta$，故 $B_{H(G)}(1,\beta)=c_{H(G)}^*(1,\beta)=k$；
\item 若 $\Phi\notin\textsc{3-SAT}$，则 $\chif(G)\ge k+1=(2+\beta)+1$（即 $\delta\ge 1$），由式 \eqref{eq:explicit_gap} 得
\[
 B_{H(G)}(1,\beta) \le k - \frac{1}{2mn} =: \tau < k = c_{H(G)}^*(1,\beta).
\]
\end{itemize}
因此 $\Phi\in\textsc{3-SAT}\iff B_{H(G)}(1,\beta)=c_{H(G)}^*(1,\beta)\iff B_{H(G)}(1,\beta)\ge k$，且 $\Phi\notin\textsc{3-SAT}\iff B_{H(G)}(1,\beta)\le\tau$，直接给出 (a), (b) 的确定性 many-one 归约与 (c) 中逆多项式加性误差 $\frac{k-\tau}{2}=\frac{1}{4mn}$ 的近似困难性。

\textbf{3. 证明 (d)：完全图 $K_K$ 的非支撑边填充与错误互斥分布可行性.}
取 $K=2n+1$ 个分类器，设定每个分类器的边际错误率为 $q_i=q:=\frac{1}{K}$（满足 $\sum_{i=1}^K q_i=1$，位于 $R_0$ 内）。取缩放因子 $\gamma:=\frac{1}{2kK}>0$，令完全图 $K_K$ 上的成对上界 $b_{uv}$ 为：
\begin{itemize}
\item 对垂直匹配边 $uv\in E_0$：$b_{uv}=0$（或严格正有理数 $\eta=\frac{\gamma}{6mn^2}$）；
\item 对星边 $uv\in E_{\rm star}$：$b_{uv}=\gamma$；对顶层边 $uv\in E_{\rm top}$：$b_{uv}=\beta\gamma=(k-2)\gamma$；
\item 对所有非支撑填充边 $uv\notin E(H(G))$：设为不起作用的平凡边际上界 $b_{uv}:=q=\frac{1}{K}$。
\end{itemize}
首先核验可行分布集 $\mathcal P(q,b)\neq\varnothing$：由于 $\sum_{i=1}^K q_i=1$，在 $[0,1]$ 上取互不相交的出错区间 $\bigl[\frac{i-1}{K},\frac{i}{K}\bigr)$（$i=1,\dots,K$）构造\textbf{错误互斥分布} $\pi$，则对一切 $u\neq v$ 均有 $\Pr_\pi(X_u=1,X_v=1)=0\le b_{uv}$，故 $\pi\in\mathcal P(q,b)\neq\varnothing$。

其次核验高成本边安全剔除条件：在支撑边上已有奇环上界 $U=k\gamma=\frac{1}{2K}$，而所有非支撑填充边的成本均为 $b_{uv}=q=\frac{1}{K}> \frac{1}{2K}=U\ge \bgraph(b)$。由于严格不等式 $b_{uv}>U$ 成立，由第一阶段报告定理 6.1（高成本边安全剔除定理），所有非支撑填充边在最优证书中的权重必为 $0$，从而 $\bgraph(b)=\gamma B_{H(G)}(1,\beta)$，且因 $\bgraph(b)\le \frac{1}{2K}<q_{\min}=\frac{1}{K}$，有 $R_{\rm det}^*(q,b)=\min\{q_{\min},\bgraph(b)\}=\gamma B_{H(G)}(1,\beta)$。证毕。
\end{proofcn}

\subsection{为什么单分支区域精确属于 $\mathrm{coNP}$：Carath\'eodory 对偶见证定理}
\begin{theorem}[单分支区域的 $\mathrm{coNP}$ 归属与对偶见证定理]\label{thm:conp_barrier}
在锚定星—顶图 $H(G)$ 上定义纯 $A$ 侧线性规划值 $B_{\text{\rm Pure-}A}(G;b)=\min\{b^\top y:\ y=(a,w)\ge 0,\ A_X(a,w)\ge 1\ (\forall X\subseteq V_G)\}$。对任意有理输入 $(G,b,\tau)$，判定“$B_{\text{\rm Pure-}A}(G;b)\le \tau$”\textbf{恒属于复杂性类 $\mathrm{coNP}$}（其补问题属于 $\mathrm{NP}$ 且为 $\mathrm{NP}$-complete）。因此，除非 $\mathrm{NP}=\mathrm{coNP}$，任何仅激活单侧分支 $A_X\ge 1$ 的图族都不可能将 $\mathrm{NP}$-complete 问题正向 Karp 归约到 $B_H(b)\le\tau$。
\end{theorem}
\begin{proofcn}
由线性规划强对偶定理，$B_{\text{\rm Pure-}A}(G;b) = \max\{\sum_{X\subseteq V_G}\lambda_X:\ \lambda_X\ge 0,\ \sum_{X\subseteq V_G}\lambda_X \nabla_y A_X\le b\}$。对偶问题仅有 $N_e=n+m$ 条线性不等式约束。由 Carath\'eodory 定理，若对偶最优值 $>\tau$，则存在至多 $n+m$ 个子集 $X_1,\dots,X_k\subseteq V_G$（$k\le n+m$）及有理权重 $\lambda_1,\dots,\lambda_k\ge 0$ 构成多项式时间可核验的 $\mathrm{NP}$ 见证。故“$B_{\text{\rm Pure-}A}(G;b)>\tau$”属于 $\mathrm{NP}$，即“$B_{\text{\rm Pure-}A}(G;b)\le\tau$”属于 $\mathrm{coNP}$。
\end{proofcn}

\section{五分量奇-$K_5$ 提升的 $\FF_2$ 上同调分类与 9 顶点 $J_4^\circ$ 唯一性}

\begin{theorem}[五分量奇-$K_5$ 提升的 $\FF_2$ 上同调分类与 9 顶点 $J_4^\circ$ 的唯一性]\label{thm:odd_K5_classification}
设 $H$ 是由 5 个零成本二部连通分量 $C_0,C_1,C_2,C_3,C_4$（每个分量 $C_u$ 为单顶点或一条零成本边 $T_uB_u$）以及连接这 5 个分量的 10 条单位成本边 $e_{uv}$（$0\le u<v\le 4$）构成的图。
\begin{enumerate}[label=(\alph*)]
\item \textbf{奇-$K_5$ 端口分裂的 $\FF_2$ 上同调刻画}：$K_5$ 的全部 10 个三角形在 $H$ 中均保持为奇环，当且仅当在适当交换各分量内部标签 $T_u\leftrightarrow B_u$ 后，存在边集二划分 $E(K_5)=E_{\rm top}\dot\cup E_{\rm bot}$，使得 $e_{uv}=T_uT_v$（当 $uv\in E_{\rm top}$）且 $e_{uv}=B_uB_v$（当 $uv\in E_{\rm bot}$）。此时分裂分量数为 $m_{\rm split}=|\{u\in V(K_5):\ \delta_{K_5}(u)\cap E_{\rm top}\neq\varnothing\ \land\ \delta_{K_5}(u)\cap E_{\rm bot}\neq\varnothing\}|$，总顶点数为 $|V(H)|=5+m_{\rm split}$。
\item \textbf{1024 种奇-$K_5$ 提升的完全分类}：在 $E(K_5)$ 的全部 $2^{10}=1024$ 种二划分中，\textbf{当且仅当 $(E_{\rm top},E_{\rm bot})$（或交换后）构成“4-星 $K_{1,4}$ 与 4-团 $K_4$”时（共 10 种，全部同构于 9 顶点图 $J_4^\circ$，$m_{\rm split}=4,\,|V(H)|=9$），有 $B_H(\1)=\frac{14}{5}<3$}；其余全部 $1014$ 种二划分（包括所有 $|V(H)|\in\{5,7,8,10\}$ 的提升）均满足 $B_H(\1)=3=c_H^*(\1)$。
\end{enumerate}
\end{theorem}
\begin{proofcn}
(a) 记边 $e_{uv}$ 在 $C_u,C_v$ 端的端口奇偶号为 $\sigma_u(uv),\sigma_v(vu)\in\FF_2$，并记 $\tau(uv):=\sigma_u(uv)\oplus\sigma_v(vu)\in\FF_2$。三角形 $(u,v,w)$ 在 $H$ 中穿过偶数条零成本边当且仅当 $(d\tau)(u,v,w)=\tau(uv)\oplus\tau(vw)\oplus\tau(wu)=0$。由 $H^1(K_5;\FF_2)=0$，存在 $\chi:V(K_5)\to\FF_2$ 使 $\tau(uv)=\chi(u)\oplus\chi(v)$；交换 $\chi(u)=1$ 处的 $T_u\leftrightarrow B_u$ 即得 $\sigma_u(uv)=\sigma_v(vu)=:\sigma(uv)\in\{0,1\}$。
(b) 由附录 A 的精确 MILP 对全部 1024 种二划分逐一求解直接验证。
\end{proofcn}

\newpage
\appendix
\section{完整可复现 Python 精确验证代码}

\begin{lstlisting}[style=pycode,caption={一般 $B_{\rm graph}$ 核心定理与 9 分类器 $R_0$ 反例精确复核脚本（\texttt{Bgraph\_breakthrough\_verify.py}）}]
#!/usr/bin/env python3
"""
B_graph Complexity & Structural Breakthroughs -- Consolidated Exact Verification Suite
Verifies:
1. [Theorem 2.1 & Corollary 2.2] Exact global optimality of B_{J_4^\circ} = 14/5 across all 2048 LP branches,
   and the 9-classifier R_0 instance with R^*_det = 7/90 < 1/12 = min(q_min, c^*).
2. [Theorem 3.1] Single-Center Anchoring Theorem over all 2^16 = 65,536 partitions (g_H = 2, ratio = 14/5).
3. [Theorem 4.3] Exact Fractional Chromatic Number Threshold Theorem on H(G):
   B_{H(G)}(1, beta) < 2 + beta  <===>  \chi_f(G) > 2 + beta.
4. [Theorem 5.1] Exhaustive classification of all 2^10 = 1024 odd-K_5 lifts:
   ONLY the 10 isomorphic copies of 9-vertex J_4^\circ (K = 9) achieve B_H < 3.
"""
from fractions import Fraction
from itertools import combinations, product
import numpy as np
import scipy.optimize as opt


def verify_task_A_J4_and_K9_voting():
    print("=" * 72)
    print("1. [Theorem 2.1 & Cor 2.2] J_4^\circ Global Optimality & 9-Classifier R_0 Instance")
    print("=" * 72)
    m = 4
    clique_pairs = list(combinations(range(m), 2))
    states = list(product([0, 1], repeat=m))
    flex_indices = [idx for idx, x in enumerate(states) if sum(x) >= 2]

    c = np.ones(10)
    best_val = 100.0
    for mask in range(1 << len(flex_indices)):
        A_ub, b_ub = [], []
        for idx, x in enumerate(states):
            use_D = ((mask >> flex_indices.index(idx)) & 1) if idx in flex_indices else 0
            row = np.zeros(10)
            if not use_D:
                for i in range(4):
                    if x[i] == 1:
                        row[i] = -1.0
                for p_idx, (i, j) in enumerate(clique_pairs):
                    if x[i] == 0 and x[j] == 0:
                        row[4 + p_idx] = -1.0
            else:
                for p_idx, (i, j) in enumerate(clique_pairs):
                    if x[i] == 1 and x[j] == 1:
                        row[4 + p_idx] = -1.0
            A_ub.append(row)
            b_ub.append(-1.0)
        res = opt.linprog(c, A_ub=A_ub, b_ub=b_ub, bounds=(0, None), method="highs")
        if res.success and res.fun < best_val:
            best_val = res.fun
    print(f"  Exhaustive min over all 2^11 = 2048 branches of J_4^\circ: {best_val:.10f} (14/5 = 2.8)")
    assert abs(best_val - 2.8) < 1e-9

    # Verify 9-classifier R_0 instance: q_i = 1/9, gamma = 1/36, filler edges = 1/9
    q_min = Fraction(1, 9)
    gamma = Fraction(1, 36)
    c_star = 3 * gamma  # 1/12 = 7.5/90
    R_det = Fraction(14, 5) * gamma  # 7/90
    print(f"  9-Classifier R_0 instance: R^*_det = {R_det} < min(q_min, c^*) = {min(q_min, c_star)}")
    assert R_det == Fraction(7, 90) and min(q_min, c_star) == Fraction(1, 12) and R_det < min(q_min, c_star)


def verify_task_B_shared_center():
    print("\n" + "=" * 72)
    print("2. [Theorem 3.1] Single-Center Anchoring (Two J_4 Copies Sharing B_0)")
    print("=" * 72)
    a, b = Fraction(2, 5), Fraction(1, 5)
    edges = []
    for offset in (0, 4):
        for i, j in combinations(range(offset + 1, offset + 5), 2):
            edges.append((("T", i), ("T", j), b, Fraction(1)))
        for i in range(offset + 1, offset + 5):
            edges.append((("B", 0), ("B", i), a, Fraction(1)))
    for i in range(1, 9):
        edges.append((("T", i), ("B", i), Fraction(2), Fraction(0)))

    var_verts = [("T", i) for i in range(1, 9)] + [("B", i) for i in range(1, 9)]
    min_g = None
    for bits in product((0, 1), repeat=16):
        color = {("B", 0): 1}
        color.update(zip(var_verts, bits))
        inner0, inner1 = Fraction(0), Fraction(0)
        for u, v, w, _ in edges:
            if color[u] == color[v]:
                if color[u] == 0:
                    inner0 += w
                else:
                    inner1 += w
        val = max(inner0, inner1)
        if min_g is None or val < min_g:
            min_g = val
    paid_cost = sum(w * cost for _, _, w, cost in edges)
    print(f"  Exhaustive check over 2^16 = 65,536 partitions: g_H = {min_g}, cost = {paid_cost}, ratio = {paid_cost/min_g}")
    assert min_g == Fraction(2) and paid_cost / min_g == Fraction(14, 5)


def compute_chi_f(n, edges_G):
    adj = [set() for _ in range(n)]
    for u, v in edges_G:
        adj[u].add(v)
        adj[v].add(u)
    A_ub = []
    for x in product([0, 1], repeat=n):
        ones = [i for i in range(n) if x[i] == 1]
        if not ones:
            continue
        if not any(v in adj[u] for i, u in enumerate(ones) for v in ones[i + 1 :]):
            row = np.zeros(n)
            for i in ones:
                row[i] = 1.0
            A_ub.append(row)
    res = opt.linprog(-np.ones(n), A_ub=A_ub, b_ub=np.ones(len(A_ub)), bounds=(0, None), method="highs")
    return -res.fun


def compute_B_HG_exact(n, edges_G, beta=1.0):
    states = list(product([0, 1], repeat=n))
    m_G = len(edges_G)
    num_e = n + m_G
    adj = [set() for _ in range(n)]
    for u, v in edges_G:
        adj[u].add(v)
        adj[v].add(u)
    flex_states, fixed_A_states = [], []
    for x in states:
        ones = [i for i in range(n) if x[i] == 1]
        if any(v in adj[u] for i, u in enumerate(ones) for v in ones[i + 1 :]):
            flex_states.append(x)
        else:
            fixed_A_states.append(x)
    total_vars = num_e + len(flex_states)
    c = np.zeros(total_vars)
    c[:n] = 1.0
    c[n:num_e] = beta
    integrality = np.zeros(total_vars)
    integrality[num_e:] = 1
    A_ub, b_ub = [], []
    for x in fixed_A_states:
        rowA = np.zeros(total_vars)
        for i in range(n):
            if x[i] == 1:
                rowA[i] = -1.0
        for e_idx, (u, v) in enumerate(edges_G):
            if x[u] == 0 and x[v] == 0:
                rowA[n + e_idx] = -1.0
        A_ub.append(rowA)
        b_ub.append(-1.0)
    for z_idx, x in enumerate(flex_states):
        z_col = num_e + z_idx
        rowA, rowD = np.zeros(total_vars), np.zeros(total_vars)
        for i in range(n):
            if x[i] == 1:
                rowA[i] = -1.0
        for e_idx, (u, v) in enumerate(edges_G):
            if x[u] == 0 and x[v] == 0:
                rowA[n + e_idx] = -1.0
            if x[u] == 1 and x[v] == 1:
                rowD[n + e_idx] = -1.0
        rowA[z_col] = -1.0
        rowD[z_col] = 1.0
        A_ub.append(rowA)
        b_ub.append(-1.0)
        A_ub.append(rowD)
        b_ub.append(0.0)
    res = opt.milp(
        c=c,
        integrality=integrality,
        bounds=opt.Bounds(np.zeros(total_vars), np.ones(total_vars)),
        constraints=opt.LinearConstraint(np.array(A_ub), -np.inf, np.array(b_ub)),
    )
    return res.fun


def verify_task_C_fractional_chromatic_equivalence():
    print("\n" + "=" * 72)
    print("3. [Theorem 4.3] Fractional Chromatic Number Threshold Theorem on H(G)")
    print("=" * 72)
    graphs = [
        ("C_5", 5, [(0, 1), (1, 2), (2, 3), (3, 4), (4, 0)]),
        ("W_5 (5-wheel)", 6, [(0, i) for i in range(1, 6)] + [(1, 2), (2, 3), (3, 4), (4, 5), (5, 1)]),
        ("K_4", 4, list(combinations(range(4), 2))),
        ("K_3 x K_2", 6, [(0, 1), (1, 2), (2, 0), (3, 4), (4, 5), (5, 3), (0, 3), (1, 4), (2, 5)]),
        ("Moser spindle", 7, [(0, 1), (0, 2), (1, 2), (1, 3), (2, 3), (0, 4), (0, 5), (4, 5), (4, 6), (5, 6), (3, 6)]),
    ]
    for name, n, edges in graphs:
        chi_f = compute_chi_f(n, edges)
        for beta in [1.0, 1.25, 1.5, 2.0]:
            b_val = compute_B_HG_exact(n, edges, beta=beta)
            thresh = 2.0 + beta
            assert (b_val < thresh - 1e-6) == (chi_f > thresh + 1e-6)
        print(f"  Verified {name:16s} (chi_f = {chi_f:.2f}): phase transition at 2 + beta == chi_f(G) holds!")


def verify_task_4_odd_K5_uniqueness():
    print("\n" + "=" * 72)
    print("4. [Theorem 5.1] Exhaustive Classification of All 1024 Odd-K_5 Lifts")
    print("=" * 72)
    k5_edges = list(combinations(range(5), 2))
    states = list(product([0, 1], repeat=4))
    strict_masks = []
    for mask in range(1024):
        sigma = [(mask >> idx) & 1 for idx in range(10)]
        has_top = [False] * 5
        has_bot = [False] * 5
        for idx, (u, v) in enumerate(k5_edges):
            if sigma[idx] == 0:
                has_top[u] = has_top[v] = True
            else:
                has_bot[u] = has_bot[v] = True
        m_split = sum(1 for u in range(5) if has_top[u] and has_bot[u])
        K = 5 + m_split

        total_vars = 10 + 16
        c_obj = np.zeros(total_vars)
        c_obj[:10] = 1.0
        integrality = np.zeros(total_vars)
        integrality[10:] = 1
        A_ub, b_ub = [], []
        for s_idx, bits in enumerate(states):
            s = (1,) + bits
            z_col = 10 + s_idx
            rowA, rowD = np.zeros(total_vars), np.zeros(total_vars)
            for e_idx, (u, v) in enumerate(k5_edges):
                if sigma[e_idx] == 1:
                    if s[u] == 1 and s[v] == 1:
                        rowA[e_idx] = -1.0
                    elif s[u] == 0 and s[v] == 0:
                        rowD[e_idx] = -1.0
                else:
                    if s[u] == 0 and s[v] == 0:
                        rowA[e_idx] = -1.0
                    elif s[u] == 1 and s[v] == 1:
                        rowD[e_idx] = -1.0
            rowA[z_col] = -1.0
            rowD[z_col] = 1.0
            A_ub.append(rowA)
            b_ub.append(-1.0)
            A_ub.append(rowD)
            b_ub.append(0.0)
        res = opt.milp(
            c=c_obj,
            integrality=integrality,
            bounds=opt.Bounds(np.zeros(total_vars), np.ones(total_vars)),
            constraints=opt.LinearConstraint(np.array(A_ub), -np.inf, np.array(b_ub)),
        )
        if res.fun < 3.0 - 1e-6:
            strict_masks.append((mask, K, res.fun))
    print(f"  Out of 1024 odd-K_5 lifts, {len(strict_masks)} have B_H < 3 (all at K = 9, B_H = 14/5 = 2.8)!")
    assert len(strict_masks) == 10 and all(K == 9 and abs(val - 2.8) < 1e-6 for _, K, val in strict_masks)


if __name__ == "__main__":
    verify_task_A_J4_and_K9_voting()
    verify_task_B_shared_center()
    verify_task_C_fractional_chromatic_equivalence()
    verify_task_4_odd_K5_uniqueness()
\end{lstlisting}

\begin{thebibliography}{99}
\bibitem{GLS1981}
M. Gr\"otschel, L. Lov\'asz, and A. Schrijver.
\newblock The Ellipsoid Method and its Consequences in Combinatorial Optimization.
\newblock \emph{Combinatorica}, 1(2):169--197, 1981.
\newblock \url{https://doi.org/10.1007/BF02579273}.

\bibitem{GLS1988}
M. Gr\"otschel, L. Lov\'asz, and A. Schrijver.
\newblock \emph{Geometric Algorithms and Combinatorial Optimization}.
\newblock Springer, 1988.
\newblock \url{https://doi.org/10.1007/978-3-642-97881-4}.

\bibitem{LundYannakakis1994}
C. Lund and M. Yannakakis.
\newblock On the Hardness of Approximating Minimization Problems.
\newblock \emph{Journal of the ACM}, 41(5):960--981, 1994.
\newblock \url{https://doi.org/10.1145/185675.306789}.

\bibitem{Zuckerman2007}
D. Zuckerman.
\newblock Linear Degree Extractors and the Inapproximability of Max Clique and Chromatic Number.
\newblock \emph{Theory of Computing}, 3(1):103--128, 2007.
\newblock \url{https://doi.org/10.4086/toc.2007.v003a006}.

\bibitem{Guenin2001}
B. Guenin.
\newblock A Characterization of Weakly Bipartite Graphs.
\newblock \emph{Journal of Combinatorial Theory, Series B}, 83(1):112--168, 2001.
\newblock \url{https://doi.org/10.1006/jctb.2001.2051}.

\bibitem{Schrijver2002}
A. Schrijver.
\newblock A Short Proof of Guenin's Characterization of Weakly Bipartite Graphs.
\newblock \emph{Journal of Combinatorial Theory, Series B}, 85(2):255--260, 2002.
\newblock \url{https://doi.org/10.1006/jctb.2001.2101}.

\bibitem{BollobasScott2002}
B. Bollob\'as and A. D. Scott.
\newblock Problems and Results on Judicious Partitions.
\newblock \emph{Random Structures \& Algorithms}, 21(3--4):414--430, 2002.
\newblock \url{https://doi.org/10.1002/rsa.10062}.

\bibitem{ShahrokhiSzekely1994}
F. Shahrokhi and L. A. Sz\'ekely.
\newblock The Complexity of the Bottleneck Graph Bipartition Problem.
\newblock \emph{Journal of Combinatorial Mathematics and Combinatorial Computing}, 15:221--226, 1994.
\newblock \url{https://combinatorialpress.com/jcmcc-articles/volume-015/the-complexity-of-the-bottleneck-graph-bipartition-problem/}.
\end{thebibliography}

\end{document}
